\section{Imersões de Sobolev}

\begin{frame}{Imersões de Sobolev}
	
\begin{defin}
	Sejam \(E\) e \(F\) espaços vetoriais normados. Dizemos que \(E\) está  \dt{imerso continuamente} em \(F\), quando existe uma aplicação linear {\color{red}injetora} \(T:E\longrightarrow F\), que também é contínua, isto é, existe \(C>0\) tal que
	\[\n{T(u)}_F\leq C \n{u}_E.\]
\end{defin}	

No caso particular das imersões de Sobolev, {\color{blue}\(T\) será sempre a inclusão} entre um mesmo espaço vetorial \(X\) com normas diferentes, isto é \(E=(X,\|\cdot \|_E)\) e \(F=(X,\|\cdot\|_F)\), a qual é sempre injetora, bastando verificar apenas que 
\[\|u\|_F\leq C\|u\|_E.\]	
Neste caso, esta imersão será denotada  por {\color{blue} \(E\hookrightarrow F\)}.

\end{frame}

\begin{frame}
\begin{exampleblock}{Pergunta}
%	Se uma função \(u\in {\color{blue}W^{1,p}(\Omega)}\), então \(u\) pertence automaticamente a algum outro espaço \({\color{red}X}\)? Em outras palavras, \( {\color{blue}W^{1,p}(\Omega)}\subset {\color{red}X}\)? 
Existe algum espaço  \({\color{red}X}\) para o qual temos automaticamente que  \({\color{blue}W^{1,p}(\Omega)}\hookrightarrow {\color{red} X }\)?
\end{exampleblock}	
A resposta é sim! Mas o espaço \({\color{red}X}\) depende de quando:
\begin{equation*}
	\begin{split}
1\leq {\color{blue}p}<\dm, \\ {\color{blue}p}=\dm, \\ \dm<{\color{blue}p}\leq \infty.
	\end{split}
\end{equation*}	
\end{frame}

\subsection*{Desigualdade de Gagliardo-Nirenberg-Sobolev}
\begin{frame}{Caso \(\Omega=\rd\) e \(1\leq p<\dm\)}
	Começemos pelo caso \(1\leq p<\dm\) e \(\Omega=\rd\). Neste caso, veremos que 
	\[{\color{blue}W^{1,p}(\rd)}\hookrightarrow {\color{red} L^q(\rd)},\]
para algum \({\color{red} q}\).

\begin{block}{Argumento de Reescala}
	A fim de encontrarmos um candidato à {\color{red} \(q\) }, tome \(u\in \cc(\rd)\) e defina
	\[u_\lambda(x)=u(\lambda x).\]
	Suponha que 
	\[\n{u_\lambda}_{{\color{red} q },\rd}\leq C \n{\nabla u}_{p,\rd}.\]
\end{block}
\end{frame}

\begin{frame}{Expoente conjugado de Sobolev}
Definimos o {\color{red}expoente conjugado de Sobolev de} \({\color{blue} p }\in [1,\dm)\) por
	\[{\color{red} p^\ast }:=\frac{\dm {\color{blue} p }}{\dm - {\color{blue} p }}.\]
Note que 
\[\frac{1}{{\color{red} p^\ast }}=\frac{1}{{\color{blue} p }}-\frac{1}{\dm}\ \  \text{ e }\ \  {\color{red}  p^\ast }>{\color{blue} p }.\]


\end{frame}

\begin{frame}{Desigualdade de Gabliardo-Nirenberg-Sobolev}
\begin{teo}[{Gabliardo-Nirenberg-Sobolev}]\label{GNS}
Se \(1\leq{\color{blue}  p }< \dm\), então \(W^{1,{\color{blue} p }}(\rd)\subset L^{{\color{red} p^\ast }}(\rd)\) e existe uma constante \(C=C_{{\color{blue} p },\dm}>0\) tal que 
\begin{equation*}
\n{u}_{{\color{red} p^\ast },\rd}\leq C\n{\nabla u}_{{\color{blue} p },\rd}
\end{equation*}
\end{teo}
\begin{exampleblock}{Passos da Prova}
\begin{enumerate}
\item Provar para \(u\in C^1_c(\rd)\).
\begin{itemize}
\item Caso \(p=1\): TFC + Desigualdade de Hölder + integrações sucessivas.
\item Caso \(p>1\): Aplicar caso anterior para \(v=|u|^\gamma\), com \(\gamma=\frac{p(\dm -1)}{\dm-p}\).
\end{itemize}
\item Usar a densidade de \(\cc(\rd)\) em \(W^{1,p}(\rd)\)
\end{enumerate}
\end{exampleblock}

\end{frame}


\begin{frame}

\begin{corol}
Se \(1\leq {\color{blue} p }<\dm \), então
\[W^{1,{\color{blue} p }}(\rd)\hookrightarrow L^q(\rd), \ \forall q\in [{\color{blue} p },{\color{red} p^\ast }].\]
\end{corol}


\begin{exampleblock}{Desigualdade de Interpolação para normas \(L^p\) {\cite[Appendix B.2(h)]{evans}}}
Sejam \(\Omega\subset \rd\) um aberto e \(1\leq p\leq q\leq \infty\). Se \(u\in L^p(\Omega)\cap L^q(\Omega)\), então \(u\in L^r(\Omega)\) para todo \(r\in [p,q]\) e 
\[
\n{u}_{r,\Omega}\leq \n{u}^\theta_{p,\Omega}\n{u}_{q,\Omega}^{1-\theta},
\]
onde
\[
\frac{1}{r}=\frac{\theta}{p}+\frac{(1-\theta)}{q}
\]
\end{exampleblock}

\end{frame}

\begin{frame}{Caso \(\Omega=\R^\dm\) e \(p=\dm\)}
\begin{corol}
Se \({\color{blue} p }=\dm\), então
\[
W^{1,{\color{blue}p}}(\rd) \hookrightarrow L^{\color{red} q }(\rd),\ \forall {\color{red} q }\in [\dm,+\infty).
\]
\end{corol}

\begin{exampleblock}{}
Usaremos a seguinte desigualdade obtida na prova do Teorema \eqref{GNS}:

 \(\forall{\color{blue}  \gamma }>1 \), temos que
\[
{\color{red} \n{u}_{{\frac{{\color{blue} \gamma}\dm}{\dm-1}}}}
\leq {\color{blue} \gamma }
{\color{red} \n{u}^{{\color{blue} \frac{1}{\gamma'}}}_{p'(\gamma-1)}}\n{\nabla u}_p^{{\color{blue} \frac{1}{\gamma}}},\ \forall u\in \cc(\rd).
\]
\end{exampleblock}

\end{frame}

%%%%%%%%%%%%%%%%%%%%%%%%%%%%%%%%%%%%%%%%%%%%%%%%%%%%%%%%%%%%%%%%%%%%%%%%
\subsection*{Desigualdade Morrey}
\begin{frame}{Caso \(\Omega=\R^\dm\) e \(p>\dm\)}
\begin{teo}[Desigualdade de Morrey]
Se \({\color{blue} p }>\dm\), então existe \(C=C_{{\color{blue}p},\dm}>0\) tal que  para toda \( u\in W^{1,{\color{blue} p }}(\rd) \) vale
\[
|u(x)-u(y)|\leq C|x-y|^{{\color{red} \gamma }} \n{\nabla u}_{1,{\color{blue} p },\rd} \text{ q.t.p } x,y\in \rd.
\]
onde \({\color{red} \gamma }:=1-\dm/{\color{blue} p }\).
Consequentemente, 
\[
\n{u}_{C^{0,{\color{red} \gamma }}(\rd)}\leq C \n{u}_{1,{\color{blue}p},\rd}\, \forall u\in W^{1,{\color{blue} p }}(\rd),
\]
o que significa que
\[
W^{1,{\color{blue}p}}(\rd) \hookrightarrow C^{0,{\color{red} \gamma }}(\rd).
\]
\end{teo}

%\begin{exampleblock}{Integral sobre hipersuperfícies em \(\rd\)}
%\[
%\int_M f(x)\, dx = \int_U f(x(u)) |\det[\frac{\partial x}{\partial u_1}]|
%\]
%\end{exampleblock}
\end{frame}

%%%%%%%%%%%%%%%%%%%%%%%%%%%%%%%%%%%%%%%%%%%%%%%%%%%%%%%%%%%%%%%%%%%%%%%%%%

%%%%%%%%%%%%%%%%%%%%%%%%%%%%%%%%%%%%%%%%%%%%%%%%%%%%%%%%%%%%%%%%%%%%%%%%
\subsection*{Imersões de Sobolev Geral}
\begin{frame}{Caso \(\Omega=\R^\dm\): Imersão Geral}
\begin{teo}[Imersão de Sobolev Geral] Para todo \(k\geq 1\) e \({\color{blue} p }\in [1,+\infty]\)
\begin{align*}
& W^{k,{\color{blue} p }}(\rd)\hookrightarrow L^{\color{red} q }(\rd), \text{ com } \frac{1}{{\color{red} q }}=\frac{1}{{\color{blue} p }}-\frac{k}{\dm},  &
 \text{ se } k{\color{blue} p }<\dm,\\
%--------------------------------------
& W^{k,{\color{blue} p }}(\rd)\hookrightarrow L^{\color{red} q }(\rd),\ \forall {\color{red} q }\in [{\color{blue} p },+\infty),& \text{ se }  k{\color{blue} p }=\dm,\\
& W^{k,{\color{blue} p }}(\rd)\hookrightarrow C^{{\color{red} m },{\color{red} \gamma }}(\rd),
%\text{ com }
%{\color{red}m}:=\bigg\lfloor k-\frac{\dm}{{\color{blue} p }}\bigg\rfloor 
%\text{ e } {\color{red} \gamma }=  \operatorname{frac}\left(k-\frac{\dm}{{\color{blue} p}}\right) 
& \text{ se }  k{\color{blue} p }>\dm,
\end{align*}
onde 
\[
\begin{cases}
{\color{red}m}:= k-\frac{\dm}{{\color{blue} p}} 
\text{ e } 0< {\color{red} \gamma }<1,
& \text{ se } k-\frac{\dm}{{\color{blue} p }} \text{é inteiro},\\
{\color{red}m}:=\bigg\lfloor k-\frac{\dm}{{\color{blue} p }}\bigg\rfloor 
\text{ e } {\color{red} \gamma }=  \operatorname{frac}\left(k-\frac{\dm}{{\color{blue} p}}\right), 
& \text{ se } k-\frac{\dm}{{\color{blue} p }} \text{ não é inteiro},
\end{cases}
\]
onde \(\lfloor\cdot \rfloor\) é a parte inteira de um número e \(\operatorname{frac(\cdot)}\) é a sua parte fracionária, isto é, \(\operatorname{frac}(x)=x-\lfloor x\rfloor\). A constante de imersão depende apenas de \(k, {\color{blue} p }\) e \(\dm\).
\end{teo}


\end{frame}

%%%%%%%%%%%%%%%%%%%%%%%%%%%%%%%%%%%%%%%%%%%%%%%%%%%%%%%%%%%%%%%%%%%%%%%%%%

\begin{frame}

\begin{exer}
Provar este resultado. Basta aplicar os resultados anteriores indutivamente. Ver \cite{evans}.
\end{exer}
\subsection*{Exercício \theexer}

\begin{exe}
\begin{enumerate}
\item \(W^{1,\infty}(\rd)\hookrightarrow C^{0,1}(\rd)\). Por isso as funções de  \(W^{1,\infty}(\rd)\) são chamadas de Lipschtzianas.
\item \(W^{1,p}(\R)\hookrightarrow C^{0,\frac{1}{2}}(\R)\), para todo \(p>1\). 

\item  \(W^{1,1}(\R)\hookrightarrow L^q(\R) \), para todo \(q\geq 1\). 

\item \(H^k(\rd)\hookrightarrow C^0(\rd)\) para \(2k>\dm\). 

\item  Se \(u\in W^{k,p}(\rd)\) para todo \(k>0\), então \(u\in C^\infty(\rd)\).
\end{enumerate}
\end{exe}
\end{frame}

%%%%%%%%%%%%%%%%%%%%%%%%%%%%%%%%%%%%%%%%%%%%%%%%%%%%%%%%%%%%%%%%%%%%%%%%%%

\subsection*{Caso \texorpdfstring{\(\Omega\subset\rd\)}{Ω⊂ℝᵈ}}
\begin{frame}{O caso \(\Omega\subset\rd\)}
\begin{teo}
Seja \(\Omega\) um aberto de \(\rd\), {\color{red} com fronteira limitada \(C^1\)} ou \(\Omega={\color{red} \rd_{+} }\). Para todo \(k\geq 1\) e \({\color{blue} p }\in [1,+\infty]\)
\begin{align*}
& W^{k,{\color{blue} p }}(\Omega)\hookrightarrow L^{\color{red} q }(\Omega), \text{ com } \frac{1}{{\color{red} q }}=\frac{1}{{\color{blue} p }}-\frac{k}{\dm},  &
 \text{ se } k{\color{blue} p }<\dm,\\
%--------------------------------------
& W^{k,{\color{blue} p }}(\Omega)\hookrightarrow L^{\color{red} q }(\Omega),\ \forall {\color{red} q }\in [{\color{blue} p },+\infty),& \text{ se }  k{\color{blue} p }=\dm,\\
& W^{k,{\color{blue} p }}(\Omega)\hookrightarrow C^{{\color{red} m },{\color{red} \gamma }}(\Omega),
%\text{ com } {\color{red}m}:=\bigg\lfloor k-\frac{\dm}{{\color{blue} p }}\bigg\rfloor 
% \text{ e } {\color{red} \gamma }=  \operatorname{frac}\left(k-\frac{\dm}{{\color{blue} p}}\right)   ,   
& \text{ se }  k{\color{blue} p }>\dm,
\end{align*}
onde 
\[
\begin{cases}
{\color{red}m}:= k-\frac{\dm}{{\color{blue} p}} 
\text{ e } 0< {\color{red} \gamma }<1,
& \text{ se } k-\frac{\dm}{{\color{blue} p }} \text{é inteiro},\\
{\color{red}m}:=\bigg\lfloor k-\frac{\dm}{{\color{blue} p }}\bigg\rfloor 
\text{ e } {\color{red} \gamma }=  \operatorname{frac}\left(k-\frac{\dm}{{\color{blue} p}}\right), 
& \text{ se } k-\frac{\dm}{{\color{blue} p }} \text{ não é inteiro},
\end{cases}
\]
A constante de imersão depende apenas de \(k, {\color{blue} p }, \dm\) e \(\Omega\).
\end{teo}
\end{frame}


%%%%%%%%%%%%%%%%%%%%%%%%%%%%%%%%%%%%%%%%%%%%%%%%
\begin{frame}{Imersões de \(W_0^{1,p}(\Omega)\)}

Para \(W_0^{1,{\color{blue} p }}(\Omega)\), valem os mesmos resultados de imersão, com uma condição mais geral para \(\Omega\).

\begin{teo}
Seja \(\Omega\) um aberto de \(\rd\). Para todo \(k\geq 1\) e \({\color{blue} p }\in [1,+\infty]\)
\begin{align*}
& W_0^{k,{\color{blue} p }}(\Omega)\hookrightarrow L^{\color{red} q }(\Omega), \text{ com } \frac{1}{{\color{red} q }}=\frac{1}{{\color{blue} p }}-\frac{k}{\dm},  &
 \text{ se } k{\color{blue} p }<\dm,\\
%--------------------------------------
& W_0^{k,{\color{blue} p }}(\Omega)\hookrightarrow L^{\color{red} q }(\Omega),\ \forall {\color{red} q }\in [{\color{blue} p },+\infty),& \text{ se }  k{\color{blue} p }=\dm,\\
& W_0^{k,{\color{blue} p }}(\Omega)\hookrightarrow C^{{\color{red} m },{\color{red} \gamma }}(\Omega),
& \text{ se }  k{\color{blue} p }>\dm,
\end{align*}
onde 
\[
\begin{cases}
{\color{red}m}:= k-\frac{\dm}{{\color{blue} p}} 
\text{ e } 0< {\color{red} \gamma }<1,
& \text{ se } k-\frac{\dm}{{\color{blue} p }} \text{é inteiro},\\
{\color{red}m}:=\bigg\lfloor k-\frac{\dm}{{\color{blue} p }}\bigg\rfloor 
\text{ e } {\color{red} \gamma }=  \operatorname{frac}\left(k-\frac{\dm}{{\color{blue} p}}\right), 
& \text{ se } k-\frac{\dm}{{\color{blue} p }} \text{ não é inteiro},
\end{cases}
\]
A constante de imersão depende apenas de \(k, {\color{blue} p }, \dm\) e \(\Omega\).
\end{teo}


\end{frame}

%%%%%%%%%%%%%%%%%%%%%%%%%%%%%%%%%%%%%%%%%%%%%%%%%%%%%%%%%%%%%%%%%%%
\subsection*{Desigualdade de Poincaré}
\begin{frame}{Desigualdade de Poincaré}

\begin{teo}[Desigualdade de Poincaré]
Seja \(\Omega\) um aberto {\color{red} limitado } de \(\rd\). Então, existe \(C_{{\color{blue}p},\dm}>0\) tal que, para todo  \({\color{blue} p }\in [1,\dm)\), vale
\begin{equation*}
\n{u}_{{\color{blue} p },\Omega}\leq C_{{\color{blue}p},\dm} |\Omega|^{\frac{1}{\dm}} \n{\nabla u}_{{\color{blue} p },\Omega},\ \forall u\in W_0^{1,{\color{blue} p }}(\Omega).
\end{equation*}
\end{teo}

\begin{exer}[{\cite[Proposition 6.12]{folland}}]
Se \(\Omega\subset \rd\) um aberto {\color{red}limitado}, então \(L^q(\Omega)\hookrightarrow L^p(\Omega)\) para todo \(q>p\geq 1\). Além disso, 
\[
\n{v}_{p,\Omega}\leq |\Omega|^{\frac{1}{p}-\frac{1}{q}}\n{v}_{q,\Omega},\ \forall v\in L^q(\Omega).
\]
\end{exer}
\subsection*{Exercício \theexer}

\end{frame}

%%%%%%%%%%%%%%%%%%%%%%%%%%%%%%%%%%%%%%%%%%%%%%%%%%%%%%%%%%%%%%%%%%%

\begin{frame}


\begin{obs}
\begin{enumerate}
\item A diferença com a desigualdade de imersão anterior é que temos apenas o gradiente de \(u \) aparecendo no lado direito.
\item Não vale para domínios ilimitados. Por exemplo, dado \(\zeta\in \D(\rd)\) tal que \(\zeta\equiv 1\) em \(B(0,1)\), \(0\leq \zeta \leq 1\) e \(\zeta\equiv 0\) em \(\rd\setminus B(0,1)\), tome
\(\zeta_n(x):=\zeta\left(\frac{x}{n}\right).\)
Então
\[
\n{\zeta_n}_{p,\rd}\to +\infty \text{ e } \n{\zeta_n}_{1,p,\rd}\to 0, \text{ quando } n\to +\infty.
\]
\end{enumerate}
\end{obs}

\begin{exer}
Leia o Remark 2.3.3 de \cite{kesavan}.
\end{exer}
\subsection*{Exercício \theexer}
\end{frame}

%%%%%%%%%%%%%%%%%%%%%%%%%%%%%%%%%%%%%%%%%%%%%%%%%%%%%%%%%%%%%%%%%%%
\subsection*{Imersões Compactas}
\begin{frame}{Imersões Compactas}
\begin{defin}
Uma imersão contínua \(E\hookrightarrow F\) é dita ser {\color{blue}compacta}, quando
todo conjunto limitado em \(E\) é {\color{blue}relativamente compacto (ou precompacto)} em \(F\). Em outras palavras, se \(\mathcal{B}\) é limitado em \(E\), então o fecho de \(\mathcal{B}\) em \(F\) é compacto.  Neste caso, usaremos a notação \(E\ic F\).
\end{defin}

\begin{prop}
Uma imersão contínua \(E\hookrightarrow F\) é dita ser {\color{blue}compacta} se, e somente se,  cada sequência limitada \((u_n)\subset E\) possui uma subsequência convergente  em \(F\). 
\end{prop}

\begin{exer}
Se \(E\hookrightarrow F \ic G\) ou  \(E\ic F \hookrightarrow G\) e \(G\) é completo, então \(E\ic G\).
\end{exer}
\subsection*{Exercício \theexer}
\end{frame}

%%%%%%%%%%%%%%%%%%%%%%%%%%%%%%%%%%%%%%%%%%%%%%%%%%%%%%%%%%%%%%%%%%%

\begin{frame}{Teorema de Ascoli-Arzelà}

\begin{teo}[Ascoli-Arzelà]
Seja \(K\) um {\color{blue}espaço métrico compacto} e \(C(K)\) o conjunto das funções contínuas \(f:K\to \R\) com a norma do \(\sup\). Se um subconjunto \(\mathcal{B}\) de   \(C(K)\)  satisfaz:
\begin{enumerate}
\item \(\mathcal{B}\) é {\color{blue}limitado} em \(C(K)\);
\item \(\mathcal{B}\) é {\color{blue}uniformemente equicontínuo}, isto é, 
\[\forall\, \ep>0, \exists\, \delta>0 \text{ tal que } d(x_1,x_2)<\delta \Rightarrow |f(x_1)-f(x_2)|<\ep,\ \forall f\in \mathcal{B};\]
\end{enumerate}
então \(\mathcal{B}\) é {\color{red}relativamente compacto} em \(C(K)\).
\end{teo}
\end{frame}

%%%%%%%%%%%%%%%%%%%%%%%%%%%%%%%%%%%%%%%%%%%%%%%%%%%%%%%%%%%%%%%%%%%

\begin{frame}{Critério de Compacidade Forte em \(L^p\)}

Denotamos a {\color{blue}translação de \(f\) por \(h\)} por \(\tau_hf(x)=f(x+h)\), \(x\in \rd\) e \(h\in \rd\).



\begin{teo}[Fréchet-Kolmogorov]
Seja {\color{blue}\(\mathcal{F}\)} um subconjunto de \(L^p(\rd)\), com \(1\leq p<\infty\). Se 
\begin{enumerate}
\item  {\color{blue}\(\mathcal{F}\)} é {\color{blue}limitado} em \(L^p(\rd)\);
\item \(\lim\limits_{|h|\to 0}\|\tau_hu-u\|_{p,\rd}=0\) {\color{blue}uniformemente} em {\color{blue}\(\mathcal{F}\)}, isto é, 
\[\forall\, \ep>0, \exists\, \delta>0 \text{ tal que } |h|<\delta \Rightarrow \|\tau_hu-u\|_{p,\rd}<\ep,\ \forall u\in {\color{blue}\mathcal{F}}.\]
\end{enumerate}
Então \({\color{blue}\mathcal{F}}\big\vert_\Omega\) é {\color{red}relativamente compacto} em \(L^p(\Omega)\), para qualquer \(\Omega\subset \rd\) com {\color{red}medida finita}. 
\end{teo}

A prova deste resultado encontra-se em  \cite[Theorem 4.26]{brezis}.
\end{frame}

%%%%%%%%%%%%%%%%%%%%%%%%%%%%%%%%%%%%%%%%%%%%%%%%%%%%%%%%%%%%%%%%%%%


\begin{frame}{Imersões Compactas}

\begin{teo}[Rellich-Kondrachov]
Seja \(\Omega \subset \rd\) um aberto {\color{blue}limitado} com {\color{red}fronteira de classe \(C^1\)}. Então as seguintes imersões são compactas:
\begin{enumerate}
\item \(W^{1,{\color{blue}p}}(\Omega)\ic L^{\color{red}q}(\Omega) \), \(1\leq {\color{red}q} <{\color{red}p^\ast}\), se \({\color{blue}p}<\dm\),
\item \(W^{1,{\color{blue}\dm}}(\Omega)\ic L^{\color{red}q}(\Omega) \), \(1\leq {\color{red}q} <\infty\), se \({\color{blue}p}=\dm\),
\item \(W^{1,{\color{blue}p}}(\Omega)\ic C^{0,{\color{red}\gamma}}(\overline{\Omega}) \), \(1\leq {\color{red}q} <{\color{red}p^\ast}\), se \({\color{blue}p}>\dm\),
\end{enumerate}
\end{teo}

\begin{lemaaux}
Se \(u\in W^{1,p}(\rd)\), então 
\[\|\tau_hu-u\|_{p,\rd}\leq |h|\,\|\nabla u\|_{p,\rd}, \forall\, h\in \rd. \]
\end{lemaaux}

\end{frame}
%%%%%%%%%%%%%%%%%%%%%%%%%%%%%%%%%%%%%%%%%%%%%%%%%%%%%%%%%%%%%%%%%%%

